\clearpage
\fancyhead[L]{\small CARLINI. EFFEMERIDI ASTRONOMICHE.}
\hypertarget{native-579-lower}{}
\begin{center}\small Göttingische gelehrte Anzeigen. No. 10. Pages 89..93. 19 January 1815.\\[3pt]\rule{.72\linewidth}{.3pt}\end{center}

{\setlength{\rightskip}{0pt plus 3em}\emph{Effemeridi astronomiche di Milano per l'anno} 1815 \emph{calcolate da} \textsc{Francesco Carlini}.
\emph{Con appendice. Dalla r.\,c. stamperia di governo.} Milan, 1814. The calendar has 128 pages and the appendix 118, in octavo.\par}

The regular articles that accompany the astronomical calendar, which retains its exemplary arrangement, and that we reviewed in earlier volumes, are supplemented this time by a very valuable addition. The addition is a list of the 34 \textsc{Maskelyne} stars taken from \textsc{Piazzi}’s new catalogue, together with the precession calculated for 1800 and 1850, the proper motion, and the constants needed to calculate aberration and nutation by Mr.\ \textsc{von Zach}’s method. It thus brings together everything needed to determine the apparent position throughout the present century and even the preceding one. The appendix contains the following important papers. Latitude of the Brera Observatory from observations of circumpolar stars, by \textsc{Barnabas Oriani}. After a well-written survey of the observational methods gradually brought into use in modern practical astronomy to determine latitudes as precisely as possible, this excellent astronomer turns to the Milan Observatory in particular. He recounts the various attempts to refine its latitude. Even the most recent observations by \textsc{Carlini}, made with a 16-inch repeating circle by \textsc{Lenoir} — the same one \textsc{Mechain} had used in Barcellona and Montjoui — left an uncertainty of several seconds of arc. Since the end of 1810 the observatory has had a 3-foot circle with an upright column by \textsc{Reichenbach}. An uninterrupted series of observations was made with it over the course of one year and printed in full in the 1812 and 1813 volumes of the Milan ephemerides. We now have the first fruits of this rich store of observations, namely the complete results for Polaris, $\delta$ Cassiopeiae and $\varepsilon$ Ursae Majoris. A survey of the individual results still shows small differences from one day to the next, amounting to several seconds of arc. These may be due partly to unavoidable observational errors and partly to changes in the atmosphere that the barometer and thermometer do not register, unless one is inclined to attribute part of the differences to the design of the instrument itself. The way the circle is attached to its column has, as is well known, raised some concerns. Because of these concerns, \textsc{Oriani} will have a movable spirit level fitted to the circle itself, and it will be most interesting to see the effect of this modification.\vadjust{\penalty-10000} Polaris was observed below the pole on 159 days and above it on 143 days. It is remarkable that the daytime observations agree less well with one another and yet give a smaller zenith distance than the nighttime observations. \textsc{Oriani} therefore excluded the daytime observations entirely and derived the following from the rest
\SourceLeaf{580}
\begin{center}\small\begin{tabular}{@{}l@{\hspace{9pt}}r@{}}
Latitude of the Brera Observatory $.\;.\;.\;.$ & $49^\circ\,28'\,0''713$ \\
Declination of Polaris for 1811 $.\;.$ & $88\phantom{\circ}\,17\phantom{'}\,59.494$ \\
\end{tabular}\end{center}
\noindent If we took the mean of \emph{all} the observations, our calculation would give
\begin{center}\small\begin{tabular}{@{}l@{\hspace{9pt}}r@{}}
Latitude of the observatory $.\;.\;.\;.\;.\;.\;.$ & $49^\circ\,28'\,1''242$ \\
Declination for 1811 $.\;.\;.\;.\;.\;.\;.\;.\;.$ & $88\phantom{\circ}\,17\phantom{'}\,59.583$ \\
\end{tabular}\end{center}
\noindent \textsc{Oriani} chose that approach because he attributes the differences in the daytime observations to particular changes in the atmosphere. Yet some part of them may well be due to annual parallax in any case. If we could disregard extraneous influences that act in the same direction during the day, always or at least predominantly, we would be inclined to give some weight to the result of our calculation, a maximum parallax in declination of $0''7$. This would correspond to a parallax in right ascension of $1^{\mathrm{s}}5$ in time, a value we might be all the readier to accept because \textsc{Piazzi} found even twice as much from the Palermo observations. --- The star $\delta$ Cassiopeiae was observed at 95 upper and 54 lower culminations, and $\varepsilon$ Ursae Majoris at 26 upper and 55 lower culminations. All these observations were made at night or in twilight, so \textsc{Oriani} found it unnecessary to exclude any of them. The resulting latitude was
\begin{center}\small\begin{tabular}{@{}l@{\hspace{9pt}}r@{}}
from $\delta$ Cassiopeiae $.\;.\;.\;.\;.\;.\;.\;.\;.\;.$ & $45^\circ\,28'\,0''975$ \\
from $\varepsilon$ Ursae Majoris $.\;.\;.\;.\;.\;.\;.$ & $45\phantom{\circ}\,28\phantom{'}\,0.\,28$ \\
\end{tabular}\end{center}
\noindent and the declinations for 1811 were
\begin{center}\small\begin{tabular}{@{}l@{\hspace{9pt}}r@{}}
of $\delta$ Cassiopeiae $.\;.\;.\;.\;.\;.\;.\;.\;.\;.$ & $59^\circ\,14'\,53''415$ \\
of $\varepsilon$ Ursae Majoris $.\;.\;.\;.\;.\;.\;.$ & $56\phantom{\circ}\,59\phantom{'}\,15.\,19$ \\
\end{tabular}\end{center}
\noindent Strictly speaking, however, this close agreement does more to establish the quality of \textsc{Carlini}’s refraction table used in the calculations than to confirm the latitude. We would like to see this table become a regular feature of the ephemerides as well. At the end \textsc{Oriani} adds a list of the declinations of 3c circumpolar stars (including the three above), without giving the results of the individual observations for them. As for the latitude, all the other stars, \textsc{Oriani} assures us, also agree to within a few tenths of a second with the result from Polaris given above.

The \emph{second} paper, also by \textsc{Oriani}, determines the opposition of Mars in 1813. At that opposition \textsc{Carlini} had observed the right ascensions with the transit telescope, and \textsc{Oriani} the declinations with the circle. The mean error of \textsc{von Lindenau}’s tables of Mars was found to be $-4''8$ in longitude and $+18''8$ in latitude, geocentrically. --- Next comes a valuable series of meridian observations of the Sun made by \textsc{Angelo Cesaris} with the mural quadrant from the beginning of 1808 to the end of 1811, continuing the observations given in the 1809 and 1810 volumes. --- The next paper, by \textsc{Carlini}, gives tables for the equation of the centre of Ceres and for the reduction of that planet to the ecliptic. The former is calculated for the eccentricity $0.0784$ and the latter for the inclination $10^\circ\,37'\,40''$. The corresponding changes are also given for a change of $0.0001$ in eccentricity and of $10''$ in inclination. Similar tables for Pallas, Juno and Vesta are promised for the next volume. We must admit that we cannot find such tables altogether suitable for Pallas and Juno, since the theory of these planets will still have to take the form of \emph{variable} elements in future, and the variations are much too large for such tables to make the calculations easier. To give an idea of the size of this variation, we note \SourceLeaf{581}that the eccentricity of the orbit of Pallas at the opposition of 1803 was $=0.24554$, whereas at the opposition of 1814 it was $=0.24135$. Given the present way of calculating planetary motions, one can quite easily do without these tables as well. The meteorological observations at the Milan Observatory in 1813 close the volume. What a fortunate climate for practical astronomy, where 188 clear days are counted in a single year!
